% !TEX root = ../main.tex
\section{Notations}
\label{sec:notations}

本文使用的主要符号见表~\ref{tab:notations}。时间以周为单位，文献中按日给出的速率须先换算；比例在计算中取 $0$ 至 $1$，在结果中可写成百分数。除特别注明外，省份和家庭型下标在各状态变量中具有相同含义。

\begingroup
\setlength{\tabcolsep}{4pt}
\renewcommand{\arraystretch}{1.12}
\setlength{\LTleft}{0pt}
% Keep stretchable right glue for the longtable caption's zero-width box.
\setlength{\LTright}{0pt plus 1fill}
\begin{longtable}{@{}>{\raggedright\arraybackslash}p{0.23\linewidth}>{\raggedright\arraybackslash}p{0.62\linewidth}>{\raggedright\arraybackslash}p{0.09\linewidth}@{}}
\caption{主要符号及说明}\label{tab:notations}\\
\toprule
符号 & 含义 & 单位 \\
\midrule
\endfirsthead
\multicolumn{3}{l}{表~\thetable\quad 主要符号及说明（续）}\\
\toprule
符号 & 含义 & 单位 \\
\midrule
\endhead
\midrule
\multicolumn{3}{r}{续下页}\\
\endfoot
\bottomrule
\endlastfoot

\multicolumn{3}{@{}l}{\textbf{索引与时间}}\\*
$t$ & 连续时间，以研究起点为零点 & 周 \\
$i,j;\ J$ & 省份索引及省份总数；本文 $J=8$ & --- \\
$k;\ K$ & 家庭型索引及家庭型总数；本文按城乡分为两类，$K=2$ & --- \\
$n;\ n_{\mathrm{obs}}$ & 观测周的序号及参与当前拟合或验证的观测周数 & --- \\
$t_{\mathrm{int}},\ t_{\mathrm{end}}$ & 干预开始时刻与情景比较的共同终点 & 周 \\

\midrule
\multicolumn{3}{@{}l}{\textbf{状态变量与人口规模}}\\*
$X_{jk}(t)$ & 第 $j$ 省、第 $k$ 类家庭型中处于状态 $X$ 的人数；$X$ 为相应舱室的通用记号 & 人 \\
$S_{jk}(t)$ & 易感人数 & 人 \\
$E_{jk}(t)$ & 已感染但尚未进入发病状态的人数 & 人 \\
$I_{jk}(t)$ & 已发病、尚未住院的人数 & 人 \\
$H_{jk}(t)$ & 住院人数；在引入 $T$ 的扩展中，不含已计入 $T$ 的患者 & 人 \\
$F_{jk}(t)$ & 死亡但尚未安葬的人数，不是累计死亡人数 & 人 \\
$R_{jk}(t)$ & 已退出模型所描述传播过程的人数，包括康复者与已安葬的死者 & 人 \\
$V_{jk}(t)$ & 接种情景中已获得保护的人数，不等同于接种剂次数 & 人 \\
$T_{jk}(t)$ & 接受指定治疗的住院人数，与 $H$ 分开计数 & 人 \\
$N_j(t),\ M_k(t)$ & 第 $j$ 省与第 $k$ 类家庭型对应的人口规模；各自的计数范围须与模型分母的定义一致 & 人 \\

\midrule
\multicolumn{3}{@{}l}{\textbf{传播与病程参数}}\\*
$\lambda_{jk}(t)$ & 第 $j$ 省、第 $k$ 类家庭型中易感者的瞬时感染风险率 & 周$^{-1}$ \\
$\beta_I,\ \beta_H,$\newline $\beta_{F,k}$ & 社区、医院及第 $k$ 类家庭型葬礼接触的传播系数；不区分家庭型时记为 $\beta_F$ & 周$^{-1}$ \\
$h_{\mathrm{fac}}$ & 院内有效接触规模的折算因子；与医院传播项配套定义 & 无量纲 \\
$k(t)$ & 相对于基准传播系数的时变强度乘子，与家庭型下标 $k$ 区分 & 无量纲 \\
$\alpha$ & 潜伏状态向发病状态转移的速率参数 & 周$^{-1}$ \\
$\gamma_h,\ \gamma_d,\ \gamma_i$ & 与社区发病者入院、死亡及康复有关的病程速率参数；具体流出项还含相应分配系数 & 周$^{-1}$ \\
$\gamma_{dh},\ \gamma_{ih}$ & 与未进入指定治疗舱的住院患者死亡及康复有关的速率参数 & 周$^{-1}$ \\
$\gamma_{dt},\ \gamma_{it}$ & 与指定治疗舱内患者死亡及康复有关的速率参数 & 周$^{-1}$ \\
$\gamma_f$ & 未安葬状态向已安葬状态转移的速率参数 & 周$^{-1}$ \\
$\theta,\ \delta$ & 总体住院比例与最终死亡概率；$\delta$ 不直接等于任一报告日的粗病死率 & 无量纲 \\
$\theta_1$ & 将总体住院比例与病程速率对应起来的派生分配系数，不直接等于 $\theta$ & 无量纲 \\
$\delta_1,\ \delta_2,$\newline $\delta_{2t}$ & 分别用于社区、普通住院及指定治疗死亡流出项的派生系数；须结合竞争转移解释 & 无量纲 \\
$a$ & 计算派生分配系数时使用的病程辅助量 & 周$^{-1}$ \\
$\Delta,\ \Sigma_H,\ \Sigma_T$ & 离开社区发病、普通住院及指定治疗状态的总速率 & 周$^{-1}$ \\

\midrule
\multicolumn{3}{@{}l}{\textbf{干预与省际流动}}\\*
$\nu$ & 接种扩展中的转移速率；从接种到获得保护的假设在模型部分说明 & 周$^{-1}$ \\
$\phi$ & 入院患者中分配至指定治疗舱的比例 & 无量纲 \\
$\varepsilon$ & 治疗情景中预设的相对效果系数，作用对象由治疗参数化确定 & 无量纲 \\
$c_{\mathrm{follow}}$ & 接触者随访率，与检出比例 $\rho$ 区分 & 无量纲 \\
$W(t),\ w_{ij}(t)$ & 省际流动率矩阵及其元素；$w_{ij}$ 表示从省 $i$ 到省 $j$ 的人均流动率 & 周$^{-1}$ \\
$\mathcal{L}$ & 由省际流动率构成的净迁移算子；$\mathcal{L}X$ 表示状态 $X$ 的迁入与迁出净变化率 & 周$^{-1}$ \\
$G,\ \gamma_m$ & 归一化的省际流动相对权重矩阵与总体流动尺度；前者无量纲，后者为速率 & 见说明 \\
$m$ & 每周出省比例；用于离散时间描述时，不与连续时间流动率直接混用 & 无量纲 \\

\midrule
\multicolumn{3}{@{}l}{\textbf{观测与参数估计}}\\*
$\rho(t)$ & 感染或发病事件进入报告数据的观测比例；采用常数近似时记为 $\rho$ & 无量纲 \\
$\rho_{\max}$ & 检出比例参数化中的上限或渐近水平，不表示实际观测到的最大值 & 无量纲 \\
$t_0,\ r$ & 检出比例变化的过渡中心与变化速率；单位分别为周与周$^{-1}$ & 见说明 \\
$y_n,\ \widehat y_n$ & 第 $n$ 周的报告新增病例数与模型给出的对应预测值 & 例 \\
$C(t),\ C^{\mathrm{rep}}(t)$ & 模型累计发病数与报告累计病例数，二者不能与单周新增混用 & 例 \\
$D(t),\ D^{\mathrm{rep}}(t)$ & 模型累计死亡数与报告累计死亡数，二者采用的统计口径分别说明 & 人 \\
$\eta$ & 时变强度拟合的平滑惩罚权重；本文另用此记号，避免与感染风险率 $\lambda_{jk}$ 混淆 & 见说明 \\

\midrule
\multicolumn{3}{@{}l}{\textbf{再生数、输入风险与评价指标}}\\*
$\mathcal{F},\ \mathcal{V}$ & 在相应分析状态下线性化得到的新感染与其余转移矩阵；花体符号不指代舱室 $F,V$ & 周$^{-1}$ \\
$\mathcal{K}$ & 下一代矩阵 & 无量纲 \\
$R_0$ & 基本再生数，在规定的全易感基准状态下定义 & 无量纲 \\
$R_{\mathrm{eff}}(t)$ & 随易感状态与干预条件变化的有效再生数 & 无量纲 \\
$R_{\mathrm{local}},\ R_j$ & 输入地区的当地再生数与第 $j$ 省的再生数 & 无量纲 \\
$q$ & 指定分支过程下，一条输入传播链最终消退的概率，不是暴发概率 & 无量纲 \\
$R^2$ & 拟合的决定系数，与移出状态 $R$ 及再生数区分 & 无量纲 \\
$\mathrm{RMSE}$ & 均方根误差，与所比较的观测量同单位；本文用于周新增病例 & 例 \\
$\mathrm{MAPE}$ & 平均绝对百分比误差；百分数形式报告，零观测值的处理须另行说明 & \% \\
$\mathrm{CFR}$ & 同一报告口径下的粗病死率，不等同于个体最终死亡概率 & \% \\
$\mathrm{LHS},\ \mathrm{PRCC}$ & 拉丁超立方抽样与偏秩相关系数；前者为抽样方法，后者为无量纲统计量 & --- \\

\end{longtable}
\endgroup

状态变量中的下标省略时，须说明表示某一分组还是对各组求和。参考材料中与病程速率有关的 $g_h,g_d,g_i$ 等记号，在本文统一写为对应的 $\gamma$ 形式；流动尺度另记为 $\gamma_m$。粗体或花体的矩阵符号与同名舱室分开使用。

表中的 $y_n$ 是一周内的病例计数，因此 RMSE 以“例”计；若将输入数据写成每周平均发病速率，RMSE 才相应以“例/周”计。$C$ 与 $D$ 仅用于累计量。平滑权重 $\eta$ 的量纲随拟合目标的归一化方式确定，在参数估计部分给出。
